% ==================================================================
%  §8.3  Taylor 公式
%  源：徐森林、薛春华《数学分析》第二册，印刷页 96-102 (PDF p. 102-108)
%  底本：PaddleOCR-VL-1.6 JSON/MD 识别结果，人工逐字核校与公式精修
%  规范：全西文标点，忠实原文定理、例题、练习题编号与推导
%  注：原书此处无插图、无思考题；练习题 8.3 共 6 题（第 4 题跨页）
% ==================================================================

\section{Taylor 公式}\label{sec:8.3}

设 $f: (x_{0} - \delta, x_{0} + \delta) \to \mathbb{R}^{1}$, $y = f(x)$ 为一元函数; 如果 $f$ 在 $x_{0}$ 有直到 $m$ 阶导数, 则 $\forall x \in (x_{0} - \delta, x_{0} + \delta)$, 已证 Peano 余项的 Taylor 公式
\[
  f(x) = \sum_{k=0}^{m}\frac{f^{(k)}(x_{0})}{k!}(x - x_{0})^{k} + o((x - x_{0})^{m}).
\]
如果 $f$ 在 $(x_{0} - \delta, x_{0} + \delta)$ 内有 $m + 1$ 阶导数, 则有 Lagrange 余项的 Taylor 公式
\begin{align*}
  f(x) &= \sum_{k=0}^{m}\frac{f^{(k)}(x_{0})}{k!}(x - x_{0})^{k} + \frac{1}{(m + 1)!}f^{(m+1)}(\xi)(x - x_{0})^{m+1}, \\
  &\qquad x \in (x_{0} - \delta, x_{0} + \delta),
\end{align*}
其中 $\xi$ 介于 $x_{0}$ 与 $x$ 之间.

现考虑 $n$ 元函数的 Taylor 公式.

\begin{theorem}[$n$ 元 Lagrange 型余项的 Taylor 公式]\label{thm:8.3.1}
设 $U \subset \mathbb{R}^{n}$ 为凸区域, $f: U \to \mathbb{R}^{1}$ 具有 $m + 1$ 阶连续偏导数, $\boldsymbol{x}^{0} = (x_{1}^{0}, x_{2}^{0}, \cdots, x_{n}^{0}) \in U$, $\boldsymbol{x} = (x_{1}, x_{2}, \cdots, x_{n}) \in U$, $\exists\, \xi \in \overline{\boldsymbol{x}^{0}\boldsymbol{x}}$ ($\boldsymbol{x}^{0}$ 与 $\boldsymbol{x}$ 的连线), s.t.
\begin{align*}
  f(\boldsymbol{x}) &= f(\boldsymbol{x}^{0}) + \sum_{k=1}^{m}\frac{1}{k!}\sum_{i_{1}, i_{2}, \cdots, i_{k}=1}^{n}\frac{\partial^{k}f}{\partial x_{i_{1}}\partial x_{i_{2}}\cdots\partial x_{i_{k}}}(\boldsymbol{x}^{0})(x_{i_{1}} - x_{i_{1}}^{0})(x_{i_{2}} - x_{i_{2}}^{0})\cdots(x_{i_{k}} - x_{i_{k}}^{0}) \\
  &\quad + \frac{1}{(m + 1)!}\sum_{i_{1}, i_{2}, \cdots, i_{m+1}=1}^{n}\frac{\partial^{m+1}f}{\partial x_{i_{1}}\partial x_{i_{2}}\cdots\partial x_{i_{m+1}}}(\boldsymbol{\xi})(x_{i_{1}} - x_{i_{1}}^{0})(x_{i_{2}} - x_{i_{2}}^{0})\cdots(x_{i_{m+1}} - x_{i_{m+1}}^{0}),
\end{align*}
或
\[
  \begin{aligned}
    f(\boldsymbol{x}) &= f(\boldsymbol{x}^{0}) + \sum_{k=1}^{m}\frac{1}{k!}\Big(\sum_{i=1}^{n}(x_{i} - x_{i}^{0})\frac{\partial}{\partial x_{i}}\Big)^{k}f(\boldsymbol{x}^{0}) \\
    &\quad + \frac{1}{(m + 1)!}\Big(\sum_{i=1}^{n}(x_{i} - x_{i}^{0})\frac{\partial}{\partial x_{i}}\Big)^{m+1}f(\boldsymbol{\xi}).
  \end{aligned}
\]
\end{theorem}

\begin{proof}
令 $\varphi(t) = f((1 - t)\boldsymbol{x}^{0} + t\boldsymbol{x})$, $t \in [0, 1]$, 显然 $\varphi \in C^{m+1}([0, 1], \mathbb{R}^{1})$, 且 $\varphi(0) = f(\boldsymbol{x}^{0})$, $\varphi(1) = f(\boldsymbol{x})$,
\[
  \varphi'(t) = \sum_{i_{1}=1}^{n}\frac{\partial f}{\partial x_{i_{1}}}((1 - t)\boldsymbol{x}^{0} + t\boldsymbol{x})(x_{i_{1}} - x_{i_{1}}^{0}),
\]
\[
  \varphi''(t) = \sum_{i_{1}, i_{2}=1}^{n}\frac{\partial^{2}f}{\partial x_{i_{1}}\partial x_{i_{2}}}((1 - t)\boldsymbol{x}^{0} + t\boldsymbol{x})(x_{i_{1}} - x_{i_{1}}^{0})(x_{i_{2}} - x_{i_{2}}^{0}),
\]
\[
  \varphi^{(k)}(t) = \sum_{i_{1}, i_{2}, \cdots, i_{k}=1}^{n}\frac{\partial^{k}f}{\partial x_{i_{1}}\partial x_{i_{2}}\cdots\partial x_{i_{k}}}((1 - t)\boldsymbol{x}^{0} + t\boldsymbol{x})(x_{i_{1}} - x_{i_{1}}^{0})\cdots(x_{i_{k}} - x_{i_{k}}^{0}),
\]
\[
  k = 1, 2, \cdots, m + 1.
\]
将 $t = 0$ 代入以上公式后得
\[
  \varphi'(0) = \sum_{i_{1}=1}^{n}\frac{\partial f}{\partial x_{i_{1}}}(\boldsymbol{x}^{0})(x_{i_{1}} - x_{i_{1}}^{0}) = \Big[\sum_{i=1}^{n}(x_{i} - x_{i}^{0})\frac{\partial}{\partial x_{i}}\Big]f(\boldsymbol{x}^{0}),
\]
\begin{align*}
  \varphi''(0) &= \sum_{i_{1}, i_{2}=1}^{n}\frac{\partial^{2}f}{\partial x_{i_{1}}\partial x_{i_{2}}}(\boldsymbol{x}^{0})(x_{i_{1}} - x_{i_{1}}^{0})(x_{i_{2}} - x_{i_{2}}^{0}) \\
  &= \Big[\sum_{i=1}^{n}(x_{i} - x_{i}^{0})\frac{\partial}{\partial x_{i}}\Big]^{2}f(\boldsymbol{x}^{0}),
\end{align*}
\[
  \vdots
\]
\begin{align*}
  \varphi^{(k)}(0) &= \sum_{i_{1}, i_{2}, \cdots, i_{k}=1}^{n}\frac{\partial^{k}f}{\partial x_{i_{1}}\partial x_{i_{2}}\cdots\partial x_{i_{k}}}(\boldsymbol{x}^{0})(x_{i_{1}} - x_{i_{1}}^{0})(x_{i_{2}} - x_{i_{2}}^{0})\cdots(x_{i_{k}} - x_{i_{k}}^{0}) \\
  &= \Big[\sum_{i=1}^{n}(x_{i} - x_{i}^{0})\frac{\partial}{\partial x_{i}}\Big]^{k}f(\boldsymbol{x}^{0}),
\end{align*}
\[
  \vdots
\]
\begin{align*}
  \varphi^{(m+1)}(\theta) &= \sum_{i_{1}, i_{2}, \cdots, i_{m+1}=1}^{n}\frac{\partial^{m+1}f}{\partial x_{i_{1}}\partial x_{i_{2}}\cdots\partial x_{i_{m+1}}}((1 - \theta)\boldsymbol{x}^{0} + \theta\boldsymbol{x}) \\
  &\quad (x_{i_{1}} - x_{i_{1}}^{0})(x_{i_{2}} - x_{i_{2}}^{0})\cdots(x_{i_{m+1}} - x_{i_{m+1}}^{0}) \\
  &= \Big[\sum_{i=1}^{n}(x_{i} - x_{i}^{0})\frac{\partial}{\partial x_{i}}\Big]^{m+1}f(\boldsymbol{\xi}), \quad \theta \in (0, 1), \quad \boldsymbol{\xi} = (1 - \theta)\boldsymbol{x}^{0} + \theta\boldsymbol{x}.
\end{align*}
于是, 将一元函数的 Lagrange 余项的 Taylor 公式用到 $\varphi$ 上得到
\[
  \varphi(1) = \varphi(0) + \sum_{k=1}^{m}\frac{\varphi^{(k)}(0)}{k!}(1 - 0)^{k} + \frac{\varphi^{(m+1)}(\theta)}{(m + 1)!}(1 - 0)^{m+1}, \quad \theta \in (0, 1),
\]
即
\[
  \begin{aligned}
    f(\boldsymbol{x}) &= f(\boldsymbol{x}^{0}) + \sum_{k=1}^{m}\frac{1}{k!}\sum_{i_{1}, i_{2}, \cdots, i_{k}=1}^{n}\frac{\partial^{k}f}{\partial x_{i_{1}}\partial x_{i_{2}}\cdots\partial x_{i_{k}}}(\boldsymbol{x}^{0})(x_{i_{1}} - x_{i_{1}}^{0})(x_{i_{2}} - x_{i_{2}}^{0})\cdots(x_{i_{k}} - x_{i_{k}}^{0}) \\
    &\quad + \frac{1}{(m + 1)!}\sum_{i_{1}, i_{2}, \cdots, i_{m+1}=1}^{n}\frac{\partial^{m+1}f}{\partial x_{i_{1}}\partial x_{i_{2}}\cdots\partial x_{i_{m+1}}}(\boldsymbol{\xi})(x_{i_{1}} - x_{i_{1}}^{0})(x_{i_{2}} - x_{i_{2}}^{0})\cdots(x_{i_{m+1}} - x_{i_{m+1}}^{0}),
  \end{aligned}
\]
或者
\begin{align*}
  f(\boldsymbol{x}) &= f(\boldsymbol{x}^{0}) + \sum_{k=1}^{m}\frac{1}{k!}\Big[\sum_{i=1}^{n}(x_{i} - x_{i}^{0})\frac{\partial}{\partial x_{i}}\Big]^{k}f(\boldsymbol{x}^{0}) \\
  &\quad + \frac{1}{(m + 1)!}\Big[\sum_{i=1}^{n}(x_{i} - x_{i}^{0})\frac{\partial}{\partial x_{i}}\Big]^{m+1}f(\boldsymbol{\xi}). \qedhere
\end{align*}
\end{proof}

在应用时, 特别重要的是 Taylor 公式的前三项, 具体表达如下:
\[
  \begin{aligned}
    f(\boldsymbol{x}) &= f(\boldsymbol{x}^{0}) + \boldsymbol{J}f(\boldsymbol{x}^{0})\begin{bmatrix} x_{1} - x_{1}^{0} \\ \vdots \\ x_{n} - x_{n}^{0} \end{bmatrix} \\
    &\quad + \frac{1}{2!}(x_{1} - x_{1}^{0}, x_{2} - x_{2}^{0}, \cdots, x_{n} - x_{n}^{0})\begin{pmatrix} \frac{\partial^{2}f}{\partial x_{1}^{2}} & \cdots & \frac{\partial^{2}f}{\partial x_{1}\partial x_{n}} \\ \vdots & & \vdots \\ \frac{\partial^{2}f}{\partial x_{n}\partial x_{1}} & \cdots & \frac{\partial^{2}f}{\partial x_{n}^{2}} \end{pmatrix}_{\boldsymbol{x}^{0}}\begin{bmatrix} x_{1} - x_{1}^{0} \\ \vdots \\ x_{n} - x_{n}^{0} \end{bmatrix} + \cdots,
  \end{aligned}
\]
其中方阵
\[
  \boldsymbol{H}f(\boldsymbol{x}^{0}) = \begin{pmatrix} \frac{\partial^{2}f}{\partial x_{1}^{2}} & \cdots & \frac{\partial^{2}f}{\partial x_{1}\partial x_{n}} \\ \vdots & & \vdots \\ \frac{\partial^{2}f}{\partial x_{n}\partial x_{1}} & \cdots & \frac{\partial^{2}f}{\partial x_{n}^{2}} \end{pmatrix}_{\boldsymbol{x}^{0}}
\]
称为函数 $f$ 的 \textbf{Hessian 矩阵}, 它是一个 $n$ 阶对称方阵 (满足 $\boldsymbol{A}^{\mathrm{T}} = \boldsymbol{A}$ 的方阵 $\boldsymbol{A}$ 称为对称方阵, 即 $a_{ji} = a_{ij}$).

\begin{theorem}[$n$ 元 Peano 型余项的 Taylor 公式]\label{thm:8.3.2}
设 $f$ 在 $\boldsymbol{x}^{0} \in \mathbb{R}^{n}$ 的某开邻域 $U \subset \mathbb{R}^{n}$ 内具有 $m + 1$ 连续偏导数, 则有
\begin{align*}
  f(\boldsymbol{x}) &= f(\boldsymbol{x}^{0}) + \sum_{k=1}^{m}\frac{1}{k!}\sum_{i_{1}, i_{2}, \cdots, i_{k}=1}^{n}\frac{\partial^{k}f}{\partial x_{i_{1}}\partial x_{i_{2}}\cdots\partial x_{i_{k}}}(\boldsymbol{x}^{0})(x_{i_{1}} - x_{i_{1}}^{0})(x_{i_{2}} \\
  &\quad - x_{i_{2}}^{0})\cdots(x_{i_{k}} - x_{i_{k}}^{0}) + R_{m}(\boldsymbol{x} - \boldsymbol{x}^{0}),
\end{align*}
其中 $R_{m}(\boldsymbol{x} - \boldsymbol{x}^{0}) = O(\|\boldsymbol{x} - \boldsymbol{x}^{0}\|^{m+1})$ 或 $o(\|\boldsymbol{x} - \boldsymbol{x}^{0}\|^{m})$, $\|\boldsymbol{x} - \boldsymbol{x}^{0}\| \to 0$.
\end{theorem}

\begin{proof}
因为 $U$ 为 $\boldsymbol{x}^{0}$ 的开邻域, 故 $\exists\, \delta > 0$, s.t. 闭球 $\overline{B(\boldsymbol{x}^{0}; \delta)} \subset U$. 又因 $f$ 的所有 $m + 1$ 阶偏导数在紧致集 $\overline{B(\boldsymbol{x}^{0}; \delta)}$ 上连续. 根据连续函数的最值定理, $\exists\, M > 0$, s.t.
\[
  M > \max_{x \in B(\boldsymbol{x}^{0}; \delta)}\Big\{\Big|\frac{\partial^{m+1}f}{\partial x_{i_{1}}\partial x_{i_{2}}\cdots\partial x_{i_{m+1}}}(\boldsymbol{x})\Big| \;\Big|\; 1 \leqslant i_{1}, i_{2}, \cdots, i_{m+1} \leqslant n\Big\}.
\]
于是,
\begin{align*}
  |R_{m}(\boldsymbol{x} - \boldsymbol{x}^{0})| &= \frac{1}{(m + 1)!}\Big|\sum_{i_{1}, i_{2}, \cdots, i_{m+1}=1}^{n}\frac{\partial^{m+1}f}{\partial x_{i_{1}}\partial x_{i_{2}}\cdots\partial x_{i_{m+1}}}(\boldsymbol{\xi}) \\
  &\quad (x_{i_{1}} - x_{i_{1}}^{0})(x_{i_{2}} - x_{i_{2}}^{0})\cdots(x_{i_{m+1}} - x_{i_{m+1}}^{0})\Big| \\
  &\leqslant \frac{M}{(m + 1)!}\Big(\sum_{i=1}^{n}|x_{i} - x_{i}^{0}|\Big)^{m+1} \leqslant \frac{M}{(m + 1)!}(n\|\boldsymbol{x} - \boldsymbol{x}^{0}\|)^{m+1} \\
  &= \frac{Mn^{m+1}}{(m + 1)!}(\|\boldsymbol{x} - \boldsymbol{x}^{0}\|)^{m+1}, \\
  &\quad \Big|\frac{R_{m}(\boldsymbol{x} - \boldsymbol{x}^{0})}{(\|\boldsymbol{x} - \boldsymbol{x}^{0}\|)^{m+1}}\Big| \leqslant \frac{Mn^{m+1}}{(m + 1)!}, \quad \boldsymbol{x} \in \overline{B(\boldsymbol{x}^{0}; \delta)} \backslash \{\boldsymbol{x}^{0}\};
\end{align*}
\[
  \Big|\frac{R_{m}(\boldsymbol{x} - \boldsymbol{x}^{0})}{(\|\boldsymbol{x} - \boldsymbol{x}^{0}\|)^{m}}\Big| \leqslant \frac{Mn^{m+1}}{(m + 1)!}\|\boldsymbol{x} - \boldsymbol{x}^{0}\| \to 0, \quad \|\boldsymbol{x} - \boldsymbol{x}^{0}\| \to 0.
\]
这就证明了 $R_{m}(\boldsymbol{x} - \boldsymbol{x}^{0}) = O(\|\boldsymbol{x} - \boldsymbol{x}^{0}\|^{m+1})$ 或 $o(\|\boldsymbol{x} - \boldsymbol{x}^{0}\|^{m})$, $\|\boldsymbol{x} - \boldsymbol{x}^{0}\| \to 0$.
\end{proof}

为了与一元函数一样能用各种方法作 Taylor 展开, 我们必须将\extref{定理 4.1.1}推广到 $n$ 元情形, 即证明 $n$ 元 Peano 型 Taylor 展开的惟一性.

\begin{theorem}[$n$ 元 Peano 型余项 Taylor 展开的惟一性]\label{thm:8.3.3}
设
\begin{align*}
  f(\boldsymbol{x}) &= \sum_{i_{1}, i_{2}, \cdots, i_{k}=1}^{m}a_{i_{1}i_{2}\cdots i_{k}}(x_{i_{1}} - x_{i_{1}}^{0})(x_{i_{2}} - x_{i_{2}}^{0})\cdots(x_{i_{k}} - x_{i_{k}}^{0}) + o(\|\boldsymbol{x} - \boldsymbol{x}^{0}\|^{m}), \quad \|\boldsymbol{x} - \boldsymbol{x}^{0}\| \to 0, \\
  f(\boldsymbol{x}) &= \sum_{i_{1}, i_{2}, \cdots, i_{k}=1}^{m}b_{i_{1}i_{2}\cdots i_{k}}(x_{i_{1}} - x_{i_{1}}^{0})(x_{i_{2}} - x_{i_{2}}^{0})\cdots(x_{i_{k}} - x_{i_{k}}^{0}) + o(\|\boldsymbol{x} - \boldsymbol{x}^{0}\|^{m}), \quad \|\boldsymbol{x} - \boldsymbol{x}^{0}\| \to 0,
\end{align*}
则 $a_{i_{1}i_{2}\cdots i_{k}} = b_{i_{1}i_{2}\cdots i_{k}}$, $i = 0, 1, 2, \cdots, m$.
\end{theorem}

\begin{proof}
为了对 $n$ 便于归纳证明, 我们将上面两式用另一种方式表达:
\begin{align*}
  f(\boldsymbol{x}) &= \sum_{k=0}^{m}\sum_{i_{1}+i_{2}+\cdots+i_{n}=k}a_{i_{1}i_{2}\cdots i_{n}}(x_{1} - x_{1}^{0})^{i_{1}}(x_{2} - x_{2}^{0})^{i_{2}}\cdots(x_{n} - x_{n}^{0})^{i_{n}} \\
  &\quad + o(\|\boldsymbol{x} - \boldsymbol{x}^{0}\|^{m}), \quad \|\boldsymbol{x} - \boldsymbol{x}^{0}\| \to 0,
\end{align*}
\begin{align*}
  f(\boldsymbol{x}) &= \sum_{k=0}^{m}\sum_{i_{1}+i_{2}+\cdots+i_{n}=k}b_{i_{1}i_{2}\cdots i_{n}}(x_{1} - x_{1}^{0})^{i_{1}}(x_{2} - x_{2}^{0})^{i_{2}}\cdots(x_{n} - x_{n}^{0})^{i_{n}} \\
  &\quad + o(\|\boldsymbol{x} - \boldsymbol{x}^{0}\|^{m}), \quad \|\boldsymbol{x} - \boldsymbol{x}^{0}\| \to 0.
\end{align*}
(归纳法) 当 $n = 1$ 时, \extref{定理 4.1.1}已证明.

假设 $n - 1$ 时, 惟一性定理已成立, 则当 $n$ 时, 将上面两式相减整理即得
\[
  \lim_{\|\boldsymbol{x} - \boldsymbol{x}^{0}\| \to 0}\frac{\sum_{k=0}^{m}\sum_{i_{1}+i_{2}+\cdots+i_{n}=k}(a_{i_{1}i_{2}\cdots i_{n}} - b_{i_{1}i_{2}\cdots i_{n}})(x_{1} - x_{1}^{0})^{i_{1}}(x_{2} - x_{2}^{0})^{i_{2}}\cdots(x_{n} - x_{n}^{0})^{i_{n}}}{\|\boldsymbol{x} - \boldsymbol{x}^{0}\|^{m}} = 0.
\]
取定 $\delta_{1}, \delta_{2}, \cdots, \delta_{n-1}$, 令 $\boldsymbol{\delta} = (\delta_{1}, \delta_{2}, \cdots, \delta_{n-1}, 1)$, $\boldsymbol{x} - \boldsymbol{x}^{0} = \boldsymbol{\delta}(x_{n} - x_{n}^{0})$, 则有
\[
  \lim_{\substack{\boldsymbol{\delta}(x_{n}-x_{n}^{0}) \to 0 \\ x_{n}-x_{n}^{0} \neq 0}}\frac{\sum_{k=0}^{m}\Big[\sum_{i_{1}+i_{2}+\cdots+i_{n}=k}(a_{i_{1}i_{2}\cdots i_{n}} - b_{i_{1}i_{2}\cdots i_{n}})\delta_{1}^{i_{1}}\delta_{2}^{i_{2}}\cdots\delta_{n-1}^{i_{n-1}}\Big](x_{n} - x_{n}^{0})^{k}}{\big(\sqrt{1 + \delta_{1}^{2} + \delta_{2}^{2} + \cdots + \delta_{n-1}^{2}}\big)^{m}|x_{n} - x_{n}^{0}|^{m}} = 0.
\]
根据 $n = 1$ 时定理成立得到
\[
  \sum_{i_{1}+i_{2}+\cdots+i_{n-1}=k-i_{n}}(a_{i_{1}i_{2}\cdots i_{n}} - b_{i_{1}i_{2}\cdots i_{n}})\delta_{1}^{i_{1}}\delta_{2}^{i_{2}}\cdots\delta_{n-1}^{i_{n-1}} = 0.
\]
再由 $\delta_{1}, \delta_{2}, \cdots, \delta_{n-1}$ 任取及定理对 $n - 1$ 成立知,
\[
  a_{i_{1}i_{2}\cdots i_{n}} - b_{i_{1}i_{2}\cdots i_{n}} = 0, \quad i_{1}, i_{2}, \cdots, i_{n} = 0, 1, 2, \cdots, m. \qedhere
\]
\end{proof}

\begin{example}\label{ex:8.3.1}
求 $f(x, y) = 2x^{2} - 3xy + y^{2} - x + y - 1$ 在 $(0, 0)$ 与 $(1, -1)$ 处的 Taylor 展开式.

\begin{solution}[解法 1]
由 $f(0, 0) = -1$, $f_{x}'(0, 0) = (4x - 3y - 1)|_{(0, 0)} = -1$, $f_{y}'(0, 0) = (-3x + 2y + 1)|_{(0, 0)} = 1$, $f_{xx}''(0, 0) = 4$, $f_{xy}''(0, 0) = -3 = f_{yx}''(0, 0)$, $f_{yy}''(0, 0) = 2$ 及三阶以上的偏导数全为零, 故 $f$ 在 $(0, 0)$ 处的 Taylor 展开式为
\[
  \begin{aligned}
    f(x, y) &= 2x^{2} - 3xy + y^{2} - x + y - 1 \\
    &= -1 - (x - 0) + (y - 0) + \frac{1}{2!}[4(x - 0)^{2} + 2(-3)xy + 2y^{2}] \\
    &= -1 - x + y + 2x^{2} - 3xy + y^{2}.
  \end{aligned}
\]
\[
  f(1, -1) = 3, \quad f_{x}'(1, -1) = (4x - 3y - 1)|_{(1, -1)} = 6,
\]
\[
  f_{y}'(1, -1) = (-3x + 2y + 1)|_{(1, -1)} = -4,
\]
\[
  f_{xx}''(1, -1) = 4, \quad f_{xy}''(1, -1) = -3 = f_{yx}''(1, -1), \quad f_{yy}''(1, -1) = 2
\]
及三阶以上的偏导数全为零, 故 $f$ 在 $(1, -1)$ 处的 Taylor 展开式为
\[
  \begin{aligned}
    f(x, y) &= 3 + 6(x - 1) - 4(y + 1) + \frac{1}{2!}\big[4(x - 1)^{2} - 6(x - 1)(y + 1) + 2(y + 1)^{2}\big] \\
    &= 3 + 6(x - 1) - 4(y + 1) + 2(x - 1)^{2} - 3(x - 1)(y + 1) + (y + 1)^{2}.
  \end{aligned}
\]
\noqedhere
\end{solution}

\begin{solution}[解法 2]
在 $(0, 0)$ 点处不必凑, 只需按幂次从小到大排, 即
\[
  f(x, y) = 2x^{2} - 3xy + y^{2} - x + y - 1 = -1 - x + y + 2x^{2} - 3xy + y^{2}.
\]
在 $(1, -1)$ 点处, 从高幂次往低幂次凑! 然后再从低幂次往高幂次重新排列. 我们有
\begin{align*}
  f(x, \boldsymbol{y}) &= 2x^{2} - 3xy + y^{2} - x + y - 1 \\
  &= 2(x - 1)^{2} - 3(x - 1)(y + 1) + (y + 1)^{2} + 4x - 2 \\
  &\quad + 3x - 3y - 3 - 2y - 2 - x + y - 1 \\
  &= 2(x - 1)^{2} - 3(x - 1)(y + 1) + (y + 1)^{2} + 6x - 4y - 7 \\
  &= 3 + 6(x - 1) - 4(y + 1) + 2(x - 1)^{2} - 3(x - 1)(y + 1) + (y + 1)^{2}. \qedhere
\end{align*}
\end{solution}
\end{example}

\begin{example}\label{ex:8.3.2}
将 $\sin(x^{2} + y^{2})$ 在 $(0, 0)$ 处作 Taylor 展开到 6 次幂.

\begin{solution}[解法 1]
令 $u = x^{2} + y^{2}$, 因 $\sin u = u - \dfrac{u^{3}}{3!} + o(u^{4})$, 所以
\[
  \begin{aligned}
    \sin(x^{2} + y^{2}) &= (x^{2} + y^{2}) - \frac{1}{3!}(x^{2} + y^{2})^{3} + o((x^{2} + y^{2})^{4}) \\
    &= x^{2} + y^{2} - \frac{1}{6}(x^{6} + 3x^{4}y^{2} + 3x^{2}y^{4} + y^{6}) + o((x^{2} + y^{2})^{4}).
  \end{aligned}
\]
\noqedhere
\end{solution}

\begin{solution}[解法 2]
$f(x, y) = \sin(x^{2} + y^{2})$, $f(0, 0) = 0$,
\[
  f_{x}'(0, 0) = 2x\cos(x^{2} + y^{2})|_{(0, 0)} = 0, \quad f_{y}'(0, 0) = 2y\cos(x^{2} + y^{2})|_{(0, 0)} = 0,
\]
\[
  f_{xx}''(0, 0) = [2\cos(x^{2} + y^{2}) - 4x^{2}\sin(x^{2} + y^{2})]\big|_{(0, 0)} = 2,
\]
\[
  f_{xy}''(0, 0) = -4xy\sin(x^{2} + y^{2})\big|_{(0, 0)} = 0 = f_{yx}''(0, 0),
\]
\[
  f_{yy}''(0, 0) = [2\cos(x^{2} + y^{2}) - 4y^{2}\sin(x^{2} + y^{2})]\big|_{(0, 0)} = 2,
\]
\[
  \vdots
\]
由二元函数的 Taylor 公式得到
\begin{align*}
  \sin(x^{2} + y^{2}) &= \frac{1}{2!}(2x^{2} + 2y^{2}) + \cdots \\
  &= x^{2} + y^{2} - \frac{1}{6}(x^{6} + 3x^{4}y^{2} + 3x^{2}y^{4} + y^{6}) + o((x^{2} + y^{2})^{3}). \qedhere
\end{align*}
\end{solution}
\end{example}

\begin{example}\label{ex:8.3.3}
将 $\mathrm{e}^{x+y}$ 在 $(0, 0)$ 处作 Taylor 展开到 4 次幂.

\begin{solution}[解法 1]
\begin{align*}
  \mathrm{e}^{x+y} &= 1 + (x + y) + \frac{1}{2!}(x + y)^{2} + \frac{1}{3!}(x + y)^{3} + \frac{1}{4!}(x + y)^{4} + o((x + y)^{4}) \\
  &= 1 + (x + y) + \frac{1}{2}(x^{2} + 2xy + y^{2}) + \frac{1}{6}(x^{3} + 3x^{2}y + 3xy^{2} + y^{3}) \\
  &\quad + \frac{1}{24}(x^{4} + 4x^{3}y + 6x^{2}y^{2} + 4xy^{3} + y^{4}) + o((x^{2} + y^{2})^{2}).
\end{align*}
\noqedhere
\end{solution}

\begin{solution}[解法 2]
\begin{align*}
  \mathrm{e}^{x+y} &= \mathrm{e}^{x}\mathrm{e}^{y} \\
  &= \Big(1 + x + \frac{x^{2}}{2!} + \frac{x^{3}}{3!} + \frac{x^{4}}{4!} + o(x^{4})\Big)\Big(1 + y + \frac{y^{2}}{2!} + \frac{y^{3}}{3!} + \frac{y^{4}}{4!} + o(y^{4})\Big) \\
  &= 1 + (x + y) + \frac{1}{2}(x^{2} + 2xy + y^{2}) + \frac{1}{6}(x^{3} + 3x^{2}y + 3xy^{2} + y^{3}) \\
  &\quad + \frac{1}{24}(x^{4} + 4x^{3}y + 6x^{2}y^{2} + 4xy^{3} + y^{4}) + o((x^{2} + y^{2})^{2}).
\end{align*}
\noqedhere
\end{solution}

\begin{solution}[解法 3]
显然, $f(0, 0) = \mathrm{e}^{x+y}|_{(0, 0)} = 1$, $\dfrac{\partial^{k}f}{\partial x_{i_{1}}\partial x_{i_{2}}\cdots\partial x_{i_{k}}}(0, 0) = \mathrm{e}^{x+y}|_{(0, 0)} = 1$. 于是, 根据二元函数的 Taylor 公式有
\begin{align*}
  \mathrm{e}^{x+y} &= 1 + (x + y) + \frac{1}{2}(x^{2} + 2xy + y^{2}) + \frac{1}{6}(x^{3} + 3x^{2}y + 3xy^{2} + y^{3}) \\
  &\quad + \frac{1}{24}(x^{4} + 4x^{3}y + 6x^{2}y^{2} + 4xy^{3} + y^{4}) + o((x^{2} + y^{2})^{2}). \qedhere
\end{align*}
\end{solution}
\end{example}

\begin{example}\label{ex:8.3.4}
证明: 当 $|x|$, $|y|$ 充分小时有
\[
  \frac{\cos x}{\cos y} \approx 1 - \frac{1}{2}(x^{2} - y^{2}).
\]

\begin{proof}[证法 1]
令 $f(x, y) = \dfrac{\cos x}{\cos y}$, 则 $f(0, 0) = 1$,
\[
  f_{x}'(0, 0) = \frac{-\sin x}{\cos y}\Big|_{(0, 0)} = 0, \quad f_{y}'(0, 0) = \frac{\cos x}{\cos^{2}y}\sin y\Big|_{(0, 0)} = 0,
\]
\[
  f_{xx}''(0, 0) = \frac{-\cos x}{\cos y}\Big|_{(0, 0)} = -1, \quad f_{xy}''(0, 0) = \frac{-\sin x\sin y}{\cos^{2}y}\Big|_{(0, 0)} = 0 = f_{yx}''(0, 0),
\]
\[
  f_{yy}''(0, 0) = \cos x\frac{\cos y\cos^{2}y - \sin y \cdot 2\cos y(-\sin y)}{\cos^{4}y}\Big|_{(0, 0)} = 1.
\]
再应用二元函数的 Taylor 公式, 有
\[
  \frac{\cos x}{\cos y} \approx 1 - \frac{1}{2}(x^{2} - y^{2}).
\]
\noqedhere
\end{proof}

\begin{proof}[证法 2 (待定系数法)]
令
\[
  \begin{aligned}
    \frac{\cos x}{\cos y} &= a_{1} + a_{2}x + a_{3}y + a_{4}x^{2} + a_{5}xy + a_{6}y^{2} + o(x^{2} + y^{2}), \\
    1 - \frac{x^{2}}{2} + o(x^{2} + y^{2}) &= (a_{1} + a_{2}x + a_{3}y + a_{4}x^{2} + a_{5}xy + a_{6}y^{2})\Big(1 - \frac{y^{2}}{2} + o(x^{2} + y^{2})\Big) \\
    &\quad + o(x^{2} + y^{2}) \\
    &= a_{1} + a_{2}x + a_{3}y + a_{4}x^{2} + a_{5}xy + \Big(a_{6} - \frac{a_{1}}{2}\Big)y^{2} + o(x^{2} + y^{2}),
  \end{aligned}
\]
比较两边 $x$ 与 $y$ 同幂次的系数得到
\[
  a_{1} = 1, \quad a_{2} = a_{3} = 0, \quad a_{4} = \frac{1}{2}, \quad a_{5} = 0,
\]
\[
  a_{6} - \frac{a_{1}}{2} = 0, \quad \text{即}\ a_{6} = \frac{a_{1}}{2} = \frac{1}{2}.
\]
于是,
\[
  \frac{\cos x}{\cos y} = 1 - \frac{1}{2}x^{2} + \frac{1}{2}y^{2} + o(x^{2} + y^{2}) \approx 1 - \frac{1}{2}(x^{2} - y^{2}). \qedhere
\]
\end{proof}
\end{example}

\begin{exblock}{练习题 8.3}
  \item 将下列多项式在指定点处展开成 Taylor 多项式:

  (1) $f(x, y) = 2x^{2} - xy - y^{2} - 6x - 3y + 5$, 在点 $(1, -2)$ 处;

  (2) $f(x, y) = x^{3} + y^{3} + z^{3} - 3xyz$, 在点 $(1, 1, 1)$ 处;

  (3) $f(x, y) = \sin(x^{2} + y^{2})$, 在点 $(0, 0)$ (到 2 次幂为止);

  (4) $f(x, y) = \dfrac{x}{y}$, 在点 $(1, 1)$ (到 3 次幂为止);

  (5) $f(x, y) = \ln(1 + x + y)$, 在点 $(0, 0)$ (到 3 次幂为止).

  \item 将 $x^{y}$ 在点 $(1, 1)$ 处作 Taylor 展开到 2 次幂.

  \item 考察二次多项式
  \[
    f(x, y, z) = (x\quad y\quad z)\begin{pmatrix} A & D & F \\ D & B & E \\ F & E & C \end{pmatrix}\begin{pmatrix} x \\ y \\ z \end{pmatrix},
  \]
  试将 $f(x + \Delta x, y + \Delta y, z + \Delta z)$ 按 $\Delta x$, $\Delta y$, $\Delta z$ 的正整数幂展开.

  \item 根据 Maclaurin 公式展开函数
  \[
    f(x, y) = \sqrt{1 - x^{2} - y^{2}}
  \]
  到 4 次幂.

  \item 当 $|x|$, $|y|$, $|z|$ 很小时, 证明:
  \[
    \cos(x + y + z) - \cos x\cos y\cos z \approx -(xy + xz + yz).
  \]

  \item 将函数 $f(x, y) = \displaystyle\int_{0}^{1}(1 + x)^{t^{2}y}\,\mathrm{d}t$ 在 $(0, 0)$ 处作 Taylor 展开到 3 次幂.
\end{exblock}
